\chapter{Network Emulation and Supporting Technologies}
\label{chap:background}

\section{Chapter Overview}
\label{sec:chapter1-overview}

This chapter introduces the technical concepts and software technologies on
which the developed framework is based. The purpose is not to provide an
exhaustive description of each tool, but to identify the functions that each
technology performs in the proposed workflow. The framework combines
container-based execution, network emulation, declarative configuration,
dynamic routing, and packet inspection. These elements must cooperate while
remaining sufficiently transparent for the user to inspect and debug a
generated laboratory.

The chapter first discusses the role of experimental network environments and
the distinction between simulation and emulation. It then presents containers
and Docker as the execution substrate used by Kathar\'a. The Kathar\'a
laboratory model is described through nodes, interfaces, collision domains,
startup files, and node directories. The final sections introduce YAML,
FRRouting, Wireshark, and the role of automated validation in a reproducible
laboratory-generation process.

\section{Experimental Network Environments}
\label{sec:experimental-network-environments}

Network experimentation is required whenever a topology, protocol, service,
or configuration must be studied without modifying an operational network. A
small experiment may contain only two hosts and one router. A more complete
scenario may include several access networks, a routing core, a server
network, multiple routing protocols, and monitoring components. Reproducing
such a scenario on physical equipment is possible, but it introduces costs
and practical limitations.

Physical laboratories require dedicated devices, cabling, power, and space.
They also make repeated experiments more difficult because the initial state
must be restored before every execution. The configuration of one device may
remain from a previous test, while a replacement device may use a different
software version. These factors reduce reproducibility and make the
experimental procedure dependent on the available hardware.

Software-based environments address these limitations by representing network
nodes and links through processes, virtual machines, containers, or simulated
models. A complete topology can be stored as files, recreated on demand, and
shared with another user; container- and namespace-based emulation systems are
commonly used to support rapid and repeatable prototyping
\cite{lantz2010network,scazzariello2020kathara}. This approach is particularly appropriate for
educational and research activities, where the objective is often to inspect
the behaviour of protocols rather than to reproduce the exact performance of
physical hardware.

The use of a software environment does not remove the need for careful design.
Addresses, routes, interfaces, and services must still be mutually coherent.
Automation is therefore valuable not only because it reduces typing, but
because it makes it possible to derive several configuration artefacts from a
single description.

\section{Simulation and Emulation}
\label{sec:simulation-emulation}

The terms \emph{simulation} and \emph{emulation} are sometimes used
interchangeably, but they describe different approaches. A simulator models
the behaviour of a system through abstractions. Events such as packet
transmission, propagation delay, queueing, or protocol transitions are
represented by the simulation engine. The simulated node does not necessarily
execute the same software that would run on a real host or router.

An emulator instead provides an execution environment in which real network
stacks and applications can operate. The nodes may be implemented as virtual
machines, containers, namespaces, or other isolated processes, but they
execute real commands and exchange real packets. This gives the user access to
standard Linux tools and routing daemons.

The distinction affects the interpretation of results. Simulation is suitable
when the experiment requires controlled mathematical models, large synthetic
networks, or repeatable event timing. Emulation is useful when the experiment
focuses on configuration, interoperability, command-line behaviour, packet
formats, and software services. The framework developed in this thesis uses
emulation because its primary goal is to generate executable Kathar\'a
laboratories rather than abstract protocol models.

The use of emulation also explains why generated files remain important. A
laboratory is not only a graph of nodes and links. Each node must receive Linux
commands, files, and, where required, routing-daemon configuration. The
framework must therefore translate an abstract topology description into a
filesystem structure that can be executed by the emulator.

\section{Virtual Machines and Containers}
\label{sec:virtual-machines-containers}

Virtual machines and containers are two common mechanisms for creating
isolated execution environments. A virtual machine emulates or virtualises a
complete computer and normally runs its own operating-system kernel. This
provides strong isolation and makes it possible to run different operating
systems on the same host. The cost is a comparatively large memory footprint,
longer startup time, and additional management overhead.

Containers isolate processes while sharing the host kernel. Each container
can have its own filesystem view, process namespace, network namespace, and
resource limits \cite{docker_container,docker_overview}. Since a container does not normally boot a complete guest
kernel, it can be created and destroyed more quickly than a virtual machine.
This property is useful for network laboratories, where a topology can contain
many nodes that execute similar Linux environments.

Containers should not be interpreted as lightweight virtual machines in every
respect. Their execution model and isolation boundaries are different. For the
purposes of this thesis, however, the relevant property is that each container
can act as an independent network node. It can own virtual interfaces, an IP
routing table, processes, and configuration files. Multiple containers can
therefore represent hosts and routers on a single workstation.

Container-based emulation also supports rapid regeneration. When the source
configuration changes, the old laboratory can be stopped and a new set of
containers can be created from the updated files. This matches the intended
workflow of the framework: the YAML description is treated as the persistent
source, while the running containers are temporary instances.

\section{Docker as the Execution Substrate}
\label{sec:docker}

Docker is the container runtime used by Kathar\'a in the developed project
\cite{docker_overview}. A Docker image contains a filesystem and metadata from
which containers are created. A container is a running or stopped instance of an image
\cite{docker_container,docker_image}. Images make it possible to provide a
consistent set of tools across nodes, while containers provide isolated
runtime instances.

Two image categories are relevant to the generator. Ordinary end hosts use the
\texttt{kathara/base} image by default. Nodes that require routing functions
use \texttt{kathara/frr}. The generator may also accept an explicit image in
the YAML configuration, allowing a node to override the default selection.
This decision is implemented by considering both the declared node type and
the presence of an enabled FRRouting section.

Docker supplies the low-level execution environment, but the framework does
not directly create Docker networks or containers. That responsibility is
delegated to Kathar\'a. This separation avoids duplicating the emulator's
functionality. The framework generates standard Kathar\'a artefacts, and
Kathar\'a subsequently translates them into Docker operations.

This architecture has an additional advantage: the generated laboratory can
be inspected and started without using the Python generator again. Once the
files exist, standard Kathar\'a commands are sufficient to operate the
laboratory. The output therefore remains portable within the Kathar\'a
ecosystem.

\section{Kathar\'a}
\label{sec:kathara}

Kathar\'a is a container-based network emulation system designed for teaching,
research, protocol development, and controlled experimentation
\cite{kathara_website,scazzariello2020kathara,bonofiglio2018kathara}. It represents network devices as
containers and interconnects their interfaces through named Layer 2 collision
domains. A laboratory is stored as a directory whose files describe the
network and the configuration of each node \cite{kathara_manpage}.

\subsection{Nodes and Interfaces}
\label{subsec:kathara-nodes-interfaces}

A Kathar\'a node is identified by a name. The node can represent an end host,
a router, a firewall, a server, or another Linux-based device. Interface
indices are explicit in \texttt{lab.conf}: index 0 maps to \texttt{eth0},
index 1 to \texttt{eth1}, and so on \cite{kathara_labconf_manpage}. In this
framework, those indices are generated from the order of the interfaces in the
YAML list.

A topology entry has the following form:

\begin{lstlisting}[caption={Basic Kathar\'a interface declaration.},label={lst:kathara-interface}]
r1[0]="lan_a"
r1[1]="backbone"
\end{lstlisting}

The first line attaches \texttt{r1:eth0} to the collision domain
\texttt{lan\_a}. The second attaches \texttt{r1:eth1} to
\texttt{backbone}. Nodes that refer to the same collision-domain name are
connected at Layer 2.

The declaration does not assign an IP address. Layer 2 connectivity and Layer
3 configuration are deliberately separated. The topology file determines
which interfaces share a network segment, while startup commands configure
addresses and routes inside the containers.

\subsection{The lab.conf File}
\label{subsec:lab-conf}

The \texttt{lab.conf} file is the central topology file of a Kathar\'a
laboratory. In addition to interface declarations, it can contain laboratory
metadata and node properties such as the selected Docker image. The generator
produces entries such as:

\begin{lstlisting}[caption={Simplified generated lab.conf content.},label={lst:lab-conf-example}]
LAB_DESCRIPTION="Example routed laboratory"
LAB_VERSION=2.0
LAB_AUTHOR="Luca-Bortoluzzi"

pc_a[0]="lan_a"
pc_a[image]="kathara/base"

r1[0]="lan_a"
r1[1]="backbone"
r1[image]="kathara/frr"
\end{lstlisting}

This file is derived from the YAML node and interface declarations. The order
of interfaces in YAML determines the corresponding interface index in
\texttt{lab.conf}. Preserving this relationship is important because the
startup file assigns addresses to \texttt{eth0}, \texttt{eth1}, and the
subsequent interfaces in the same order.

\subsection{Startup Files}
\label{subsec:startup-files}

For a node named \texttt{r1}, Kathar\'a executes the commands stored in
\texttt{r1.startup} when the laboratory is started. Startup files commonly
configure addresses, default gateways, static routes, forwarding, and
services. A simplified file may contain:

\begin{lstlisting}[language=bash,caption={Simplified node startup file.},label={lst:startup-example}]
ip address add 10.10.1.1/24 dev eth0
ip address add 10.0.12.1/30 dev eth1
systemctl start frr
\end{lstlisting}

The generator writes one startup file for each node. It also appends custom
commands declared in YAML. This mechanism keeps common configuration
automatic while retaining an escape hatch for scenario-specific behaviour.

\subsection{Node Directories and hostlab}
\label{subsec:node-directories}

A directory with the same name as a node contains files that are made
available inside that node. In the project workflow, resources stored beside a
YAML configuration can be copied to the generated laboratory when the
\cliopt{import-dirs} option is used. A directory named \texttt{server}, for
example, is imported only if the YAML file declares a node named
\texttt{server}.

These resources can be accessed from the container through the laboratory
mount, commonly under \texttt{/hostlab}. This allows the startup file to
execute scripts that remain visible and editable on the host. The generator
excludes cache files and repository metadata so that generated laboratories do
not contain unnecessary development artefacts.

\subsection{Laboratory Lifecycle}
\label{subsec:kathara-lifecycle}

A generated laboratory is started with \texttt{kathara lstart} from its
directory. Kathar\'a reads \texttt{lab.conf}, creates the containers,
configures their interfaces, mounts node directories, and executes startup
files. A running node can be accessed using \texttt{kathara connect}.

The laboratory is stopped and removed from the runtime with
\texttt{kathara lclean}. The source files remain on disk and can be used for a
later execution \cite{kathara_manpage}. The project provides wrapper scripts for these operations;
their implementation and the complete generation lifecycle are discussed in
the following chapter.

\section{Declarative Configuration with YAML}
\label{sec:yaml-background}

YAML is a human-readable data-serialisation language frequently used for
configuration files \cite{yaml122}. It supports mappings, lists,
scalars, comments, and nested structures. These properties make it suitable
for representing a network laboratory in a form that is easier to maintain
than several independent startup and routing files.

A declarative description states the desired structure of the laboratory. It
does not require the user to write every command needed to create that
structure. The following fragment describes a host and a router:

\begin{lstlisting}[caption={Simplified YAML topology declaration.},label={lst:yaml-background-example}]
lab_name: example_lab

nodes:
  pc_a:
    type: host
    interfaces:
      - network: lan_a
        ip: 10.10.1.10/24
    default_gateway: 10.10.1.1

  r1:
    type: router
    interfaces:
      - network: lan_a
        ip: 10.10.1.1/24
      - network: backbone
        ip: 10.0.12.1/30
\end{lstlisting}

The same data are later translated into several outputs: interface entries in
\texttt{lab.conf}, address commands in startup files, a default route for the
host, and potentially routing configuration for the router. The YAML document
therefore becomes the persistent source of truth for generated files.

Using a structured input also enables validation. The framework can check
whether required fields exist, addresses are valid and unique, gateways belong
to a local subnet, and routing declarations refer to existing interfaces. The
validation process is a central subject of Chapter~\ref{chap:configuration-validation}.

\section{Dynamic Routing and FRRouting}
\label{sec:frrouting-background}

FRRouting, commonly abbreviated as FRR, is an open-source routing suite for
Linux and Unix systems. It provides implementations of protocols including
OSPF, RIP, and BGP \cite{frr_user_guide}. In the generated laboratories,
FRR allows containerised nodes to exchange routes dynamically and behave as
software routers.

FRR is organised around a set of protocol daemons and a common configuration
interface. The \texttt{daemons} file determines which processes are enabled.
The integrated \texttt{frr.conf} file contains protocol configuration, while
\texttt{vtysh} provides a unified command-line interface for inspection and
administration.

\subsection{OSPF}
\label{subsec:ospf-background}

Open Shortest Path First is a link-state interior gateway protocol. Routers
exchange topology information and compute paths within an autonomous routing
domain \cite{rfc2328}. The project supports both a common area declared at node level and an
\texttt{ospf\_area} marker associated with individual interfaces. This allows
simple single-area topologies and more structured multi-area configurations.

The generator can derive OSPF network statements from interface prefixes.
Validation verifies that explicitly declared OSPF networks overlap a locally
connected interface. This prevents a router from advertising a prefix that is
unrelated to its own configuration.

\subsection{RIP}
\label{subsec:rip-background}

Routing Information Protocol is a distance-vector protocol with a simpler
configuration model \cite{rfc2453}. In the project, the generator can enable the RIP daemon,
select the protocol version, disable automatic summarisation, and derive
network statements from node interfaces. RIP is useful for demonstrating that
the framework is not tied to one routing protocol.

\subsection{BGP}
\label{subsec:bgp-background}

Border Gateway Protocol exchanges reachability information between autonomous
systems \cite{rfc4271}. A BGP configuration requires an autonomous-system number, a router
identifier, one or more neighbours, and the prefixes to advertise. The
framework validates neighbour addresses, remote autonomous-system values, and
the existence of corresponding peer nodes.

The generated configuration includes two commands that are relevant in small
FRR laboratories: \texttt{no bgp ebgp-requires-policy} and
\texttt{no bgp network import-check}. Their use can be controlled through the
configuration model. They allow deliberately minimal educational topologies
to advertise prefixes without requiring additional policy definitions or an
already installed matching route. The policy-related behaviour is particularly
relevant to the default eBGP propagation model standardised in RFC~8212
\cite{rfc8212}.

\section{Traffic Inspection with Wireshark}
\label{sec:wireshark-background}

Wireshark is a network protocol analyser used to capture and inspect packets
\cite{wireshark_user_guide}. In an emulated environment it is useful for confirming
that packets traverse the expected collision domain, inspecting protocol
messages, and distinguishing routing failures from application failures.

The framework supports optional real-time Wireshark integration by adding a
dedicated node directly to \texttt{lab.conf}. The user can attach the node to
one or more selected collision domains. This is preferable to modifying an
existing host because monitoring remains separated from the logical role of
the emulated devices.

Packet capture is optional. A laboratory can be generated and executed without
Wireshark, while the same topology can later be regenerated with monitoring
enabled. This reflects a general design principle of the framework: diagnostic
features should not be mandatory components of every laboratory.

\section{Automation, Validation, and Reproducibility}
\label{sec:automation-reproducibility}

Automation is often described as a means of reducing manual effort, but its
more important contribution in this project is consistency. A manually edited
laboratory may contain a valid \texttt{lab.conf} and valid startup files that
nevertheless disagree with one another. A declarative generator can derive
both artefacts from the same interface list.

The framework applies two levels of validation before generation. Structural
validation verifies the minimum document shape required by the generator.
Semantic validation analyses relationships among nodes, interfaces, subnets,
gateways, and routing declarations. Errors stop generation, while warnings
identify suspicious but potentially intentional configurations.

Reproducibility depends on preserving the source configuration and generation
parameters. The generated laboratory is treated as a derived artefact. If a
persistent change is required, it should be made in the YAML file or in the
node resources stored under \texttt{configs/}, not only in the generated
\texttt{labs/} directory. A clean regeneration can then reconstruct the
laboratory from the maintained source.

\section{Positioning with Respect to Related Work}
\label{sec:related-work-positioning}

The framework combines ideas that are usually addressed separately in the
literature: executable network emulation, declarative configuration,
deterministic validation, and collection of application-level evidence. Its
contribution is not a new container runtime or routing suite. Rather, it is the
integration of these responsibilities around the standard Kathar\'a laboratory
format.

Mininet demonstrated that realistic network software could be used for rapid
prototyping on a single machine, with particular emphasis on software-defined
network experiments \cite{lantz2010network}. Kathar\'a provides the execution
substrate selected in this thesis and represents a broader range of Linux
network nodes through containers and collision domains
\cite{bonofiglio2018kathara,scazzariello2020kathara}. The developed framework
operates one level above Kathar\'a: it does not replace its runtime or file
format, but generates those files from a validated declarative source and adds
repeatable lifecycle and measurement operations.

A second line of work studies the use of LLMs for network configuration.
NetConfEval evaluates whether LLMs can perform network-configuration tasks,
while more recent work considers their role in zero-touch management and
surveys their wider use in network operations
\cite{wang2024netconfeval,lira2025llmnetcfg,liu2025llmnetworking}. These works
motivate natural-language assistance, but they also address a different
abstraction boundary. The present framework does not attempt to infer intent:
it receives explicit YAML and applies repeatable rules before producing an
executable laboratory. An LLM can therefore be used upstream as an optional
authoring aid, without becoming the authority that accepts a configuration.

\begin{table}[htbp]
\centering
\small
\setlength{\tabcolsep}{4pt}
\begin{tabular}{>{\raggedright\arraybackslash}p{0.16\textwidth}>{\raggedright\arraybackslash}p{0.25\textwidth}>{\raggedright\arraybackslash}p{0.25\textwidth}>{\raggedright\arraybackslash}p{0.22\textwidth}}
\toprule
\textbf{Approach} & \textbf{Scope and abstraction} & \textbf{Validation or evidence boundary} & \textbf{Relation to this work} \\
\midrule
Mininet \cite{lantz2010network} & Rapid prototyping, especially of SDN networks, through programmatic topologies and host/switch configuration & Experiment-specific correctness checks remain the responsibility of the experimenter & Alternative emulation approach; not the selected output target \\
Kathar\'a \cite{bonofiglio2018kathara,scazzariello2020kathara} & Container-based execution from laboratory files, startup scripts, and node resources & Consumes and executes its own laboratory representation & Runtime and file format used by the developed framework \\
LLM-based network configuration \cite{wang2024netconfeval,lira2025llmnetcfg,liu2025llmnetworking} & Translation of natural-language intent and assistance with operational tasks & Generated output remains candidate configuration requiring technical verification & Optional upstream authoring and diagnostic support \\
This thesis & Validated generation and repeatable evaluation from declarative YAML & Structural and semantic checks precede generation; runtime measurements test observable behaviour & Integration of source configuration, standard artefacts, lifecycle, and evidence \\
\bottomrule
\end{tabular}
\caption{Positioning of the developed framework relative to the main related approaches. The rows describe complementary scopes rather than a performance ranking.}
\label{tab:related-work-positioning}
\end{table}

The distinction in Table~\ref{tab:related-work-positioning} also limits the
claims of the thesis. The work does not compare emulator throughput, evaluate
LLM accuracy, or propose an alternative routing implementation. It evaluates
whether a narrow, explicit configuration model can coordinate existing tools,
reject recognised inconsistencies before generation, and support a controlled
experiment through a repeatable workflow.

\section{Chapter Summary}
\label{sec:chapter1-summary}

This chapter introduced the execution model and technologies used by the
framework. Container-based emulation provides isolated Linux network nodes,
Docker supplies the runtime, and Kathar\'a creates the topology from standard
laboratory files. YAML provides a declarative source representation,
FRRouting supplies dynamic routing, and Wireshark provides optional packet
inspection.

The chapter also established the distinction between source configuration and
generated artefacts. This distinction is central to the design of the system.
The next chapter translates the technological context into explicit
requirements, responsibilities, and architectural components.
